% Generated draft: requires mathematical and source review.
\documentclass[11pt]{article}
\usepackage[margin=25mm]{geometry}
\usepackage{fontspec}
\setmainfont{Latin Modern Roman}
\usepackage{amsmath,amssymb,mathtools,unicode-math}
\setmathfont{Latin Modern Math}
\usepackage{enumitem,needspace,xurl,hyperref}
\setlength{\parindent}{0pt}
\setlength{\parskip}{5pt}
\setlength{\emergencystretch}{3em}
\newcommand{\N}{\mathbb{N}}
\newcommand{\Z}{\mathbb{Z}}
\newcommand{\Q}{\mathbb{Q}}
\newcommand{\R}{\mathbb{R}}
\newcommand{\C}{\mathbb{C}}
\newcommand{\F}{\mathbb{F}}
\newcommand{\A}{\mathbb{A}}
\newcommand{\PP}{\mathbb{P}}
\newcommand{\E}{\mathbb{E}}
\newcommand{\eps}{\varepsilon}
\newcommand{\dd}{\,\mathrm{d}}
\begin{document}
\section*{UnsolvedMath ID: 30000011}
\textbf{Source problem code:} catalogue entry
\subsection{Definitions and mathematical statement}
\textbf{Question.} If a positive closed current \(T\) with potential \(\varphi\) satisfies\[T[\varphi]\ge\lambda[C]\qquad\text{for some }\lambda>0\text{ and a curve }C,\] must the geometric genus of \(C\) be less than \(2\)? For a reduced irreducible compact complex curve, the geometric genus means \(g(\widetilde C)=\dim_{\C}H^0(\widetilde C,\Omega^1_{\widetilde C})\), where \(\widetilde C\) is its smooth normalization; the desired conclusion is then \(g(\widetilde C)\in\{0,1\}\). The integration current \([C]\) acts on compatible test forms by integration along the regular part of \(C\). An inequality of real currents means that the difference is a positive current of the same degree.\par The source does not define the operation \(T[\varphi]\), the ambient complex space, the degree of \(T\), the sense of potential, or the hypotheses on \(C\). Accordingly, \(T[\varphi]\) is retained as an undefined source-specific notation; it is not identified with \(T\), a wedge product, a restriction, or a regularization. The normalization definition above is conditional on \(C\) being reduced, irreducible, and compact, properties not stated in the extracted question. A precise quantified conjecture requires the missing construction and geometric hypotheses; positivity and closedness alone do not specify them. In particular, no globally defined potential or ambient compact surface is imposed silently.\par The supplied literature triage dated 2026-08-21 labels this partially solved. It reports that the genus bound is proved in important foliated-surface cases, whereas the general positive-current formulation was described as only partially advanced and no general proof or counterexample was located in that triage. No precise foliated-surface theorem or its hypotheses are supplied here, so that report does not fill the formulation gap or certify the general question's status.
\subsection{Short English statement}
The source asks whether a curve detected with positive coefficient by the current construction \(T[\varphi]\) must have normalization of genus zero or one. The construction and ambient hypotheses are missing, so the inequality is preserved without assigning it an unsupported interpretation.
\subsection{Sources}
\begin{itemize}\item Local UnsolvedMath registry, permanent ID 30000011, code catalogue entry; \emph{Genus Bounds from Positive Closed Currents}, Oberwolfach Reports Open Problems.\item \emph{Multiplier Ideal Sheaves in Algebraic and Complex Geometry (2005)}, Oberwolfach Reports, DOI 10.4171/owr/2004/21, pp. 1108--1166. \url{https://doi.org/10.4171/owr/2004/21}. Proposed year and year tag: 2004; no proposer supplied.\item Supplied literature label \emph{original OWR report}: \url{https://doi.org/10.4171/owr/2004/21}.\item Supplied literature label \emph{McQuillan, curves on surfaces}: \url{https://www.mat.uniroma2.it/~mcquilla/files/charp.pdf}. No fuller citation or theorem statement is supplied.\item Literature triage checked 2026-08-21, status \texttt{partially\_solved}; record updated 2026-09-16. The registry's \texttt{corrected\_verified} label records restoration of notation from supplied prose, not a supplied definition of \(T[\varphi]\). Verify the dated assessment independently.\end{itemize}
\end{document}
